\documentclass[10pt]{beamer}

\usepackage{verbatim}
%\usepackage[latin1]{inputenc}
\usepackage{amssymb,amsmath,amsfonts,url}
\usepackage{graphicx,color}
\usepackage{pgf,pgfarrows,pgfnodes,pgfautomata,pgfheaps,pgfshade}
\usepackage[utf8]{inputenc}
\usepackage[french]{babel}
\usepackage{tikz}

\usepackage[all]{xy}
\usepackage{amstext}
\usepackage{xspace}
\usepackage{mdwtab}


\mode<article> {
  \usepackage{fullpage}
  \usepackage{hyperref}
  \setjobnamebeamerversion{advancedcourse.beamer}
}




\usepackage{listings}
\usepackage{xcolor}

\lstset{
	language=Python,
	basicstyle=\ttfamily\small,
	keywordstyle=\color{blue},
	commentstyle=\color{gray},
	stringstyle=\color{green!50!black},
	showstringspaces=false,
	frame=single,
	numbers=none
}
\lstset{
	language=Python,
	basicstyle=\ttfamily\small,
	keywordstyle=\color{blue},
	commentstyle=\color{gray},
	stringstyle=\color{green!50!black},
	showstringspaces=false,
	frame=single,
	numbers=none,
	literate=
	{é}{{\'e}}1
	{è}{{\`e}}1
	{ê}{{\^e}}1
	{ë}{{\"e}}1
	{à}{{\`a}}1
	{â}{{\^a}}1
	{ä}{{\"a}}1
	{ù}{{\`u}}1
	{û}{{\^u}}1
	{ü}{{\"u}}1
	{ô}{{\^o}}1
	{ö}{{\"o}}1
	{î}{{\^i}}1
	{ï}{{\"i}}1
	{ç}{{\c c}}1
}

\newcommand{\alertonce}[1]{\addtocounter{beamerpauses}{-1}\alert<+>{#1}}
\newcommand{\alerttwice}[1]{\addtocounter{beamerpauses}{-1}\alert<+-+(1)>{#1}}

%\newcommand{\subsubsubsection}[1]{\noindent\textbf{#1}}

\newcommand{\cfbox}[1]{\begin{center}\fbox{#1}\end{center}}
\newcommand{\cbox}[1]{\begin{center} #1 \end{center}}

%%%%%%%%%%%%%%%%%%%%%%%%%%%%%%%%%%%%%%%%%%%%%%%%%%%%%%%%%%%%%%%%%%%%%%%%%%%%%%%%%%%%%%%%%

\graphicspath{{Graphics/}}
%\input def_0.tex

%\input{StyleSlides.tex}
\mode<presentation> {
    \usepackage{beamerthemeBoadilla}

  %\beamertemplatetransparentcovereddynamic
  \beamertemplateballitem
}


\title[IN100 Cours 6]{\textbf{Fondements informatiques I}\\
Cours 6: Listes}

\author[]{Sorina Ionica \texttt{sorina.ionica@uvsq.fr} \\~\\Bilal FAYE  \texttt{bilal.faye@uvsq.fr} }


\institute{}

\date{}


%\setbeamercolor{normal text}{fg=black}

\begin{document}


\frame{\titlepage}

\begin{frame}[fragile]{Les listes : une collection d'éléments}

Un \textbf{tableau} est une structure de données qui regroupe
plusieurs éléments, chacun occupant une position.

\begin{exampleblock}{Exemple : les moyens de transport}
\begin{center}
\begin{tabular}{c|c|c|c|c}
\textbf{Indice} & 0 & 1 & 2 & 3 \\ \hline
\textbf{Élément} & bus & voiture & vélo & trottinette
\end{tabular}
\end{center}
\end{exampleblock}

\begin{alertblock}{À retenir}
Chaque élément possède un \textbf{indice} qui indique sa position.
En Python, le premier indice est \textbf{0}.
\end{alertblock}

\pause

En Python, on utilise le type \texttt{list} pour représenter une liste.
Les éléments sont placés entre crochets \texttt{[]} et séparés par des virgules.

\begin{exampleblock}{Créer des listes}
\begin{lstlisting}[language=Python]
# Liste de capitales européennes
capitales = ["Paris", "Berlin", "Madrid"]
# Liste d'entiers
entiers = [5, 0, 17, 42, 23]
\end{lstlisting}
\end{exampleblock}

\end{frame}


\begin{frame}[fragile]{Créer des listes en Python}

La construction des listes en Python est très flexible :
une liste peut contenir des éléments de \textbf{types différents}.

\begin{exampleblock}{Exemple : une liste mixte}
\begin{lstlisting}[language=Python]
mixte = [7, "tennis", True, 5.3]
\end{lstlisting}
\end{exampleblock}

Cette liste contient quatre valeurs de types différents :

\begin{itemize}
    \item \texttt{7} : entier (\texttt{int}) ;
    \item \texttt{"tennis"} : chaîne de caractères (\texttt{str}) ;
    \item \texttt{True} : booléen (\texttt{bool}) ;
    \item \texttt{5.3} : nombre décimal (\texttt{float}).
\end{itemize}

\begin{block}{À retenir}
Une liste Python peut regrouper des valeurs de types différents.
\end{block}

\end{frame}


\begin{frame}[fragile]{Connaître la taille d'une liste}

La fonction intégrée \texttt{len()} permet de connaître le
\textbf{nombre d'éléments} d'une liste.

\begin{exampleblock}{Exemple}
\begin{lstlisting}[language=Python]
liste = [0, 1, 2, 3, 4, 5, 6, 7, 8, 9]

print("La liste contient", len(liste), "éléments.")
\end{lstlisting}
\end{exampleblock}


\begin{block}{Résultat}
\begin{lstlisting}[language=Python]
La liste contient 10 éléments.
\end{lstlisting}
\end{block}

\begin{alertblock}{Attention}
La liste contient 10 éléments, mais son dernier indice est
\textbf{9}, car les indices commencent à 0.
\end{alertblock}

\end{frame}
\begin{frame}[fragile]{Modifier une liste}

La taille d'une liste Python est \textbf{variable} :
on peut ajouter ou supprimer des éléments après sa création.

\vspace{0.3cm}

\begin{exampleblock}{Ajouter un élément avec \texttt{append()}}
\begin{lstlisting}[language=Python]
fruits = ["pomme", "banane"]
fruits.append("orange")
print(fruits)  # ['pomme', 'banane', 'orange']
\end{lstlisting}
\end{exampleblock}

\begin{exampleblock}{Supprimer un élément avec \texttt{remove()}}
\begin{lstlisting}[language=Python]
fruits.remove("banane")
print(fruits)  # ['pomme', 'orange']
\end{lstlisting}
\end{exampleblock}

\begin{alertblock}{À retenir}
\begin{itemize}
    \item \texttt{append(element)} ajoute un élément à la fin.
    \item \texttt{remove(element)} supprime la première occurrence de cet élément.
\end{itemize}
\end{alertblock}

\end{frame}
\begin{frame}[fragile]{Accéder aux éléments d'une liste}

Chaque élément d'une liste possède un \textbf{indice} qui indique sa position.

\begin{itemize}
    \item Le premier indice est \texttt{0}.
    \item Pour une liste de longueur $n$, les indices vont de $0$ à $n-1$.
    \item Un indice doit être un nombre entier.
\end{itemize}

\begin{exampleblock}{Exemple}
\begin{lstlisting}[language=Python]
transports = ["bus", "voiture", "velo", "trotinette"]

print(transports[0])  # bus
print(transports[2])  # velo
\end{lstlisting}
\end{exampleblock}

\begin{alertblock}{Attention}
Un indice qui n'existe pas dans la liste provoque une erreur.
Pour une liste de 4 éléments, \texttt{transports[4]} est invalide.
\end{alertblock}

\end{frame}


\begin{frame}[fragile]{Accéder aux éléments depuis la fin}

Python permet aussi d'utiliser des \textbf{indices négatifs} :
on compte alors à partir de la fin de la liste.

\vspace{0.3cm}

\begin{center}
\begin{tabular}{c|c|c|c|c}
Élément & bus & voiture & vélo & trottinette \\ \hline
Indice positif & 0 & 1 & 2 & 3 \\ \hline
Indice négatif & -4 & -3 & -2 & -1
\end{tabular}
\end{center}

\begin{exampleblock}{Exemple}
\begin{lstlisting}[language=Python]
transports = ["bus", "voiture", "vélo", "trottinette"]

print(transports[-1])  # trottinette
print(transports[-2])  # vélo
print(transports[0])   # bus
print(transports[-4])  # bus
\end{lstlisting}
\end{exampleblock}

\begin{block}{Avantage}
Avec \texttt{-1} et \texttt{-2}, on accède au dernier et à
l'avant-dernier élément sans connaître la longueur de la liste.
\end{block}

\end{frame}


\begin{frame}[fragile]{Extraire une partie d'une liste : les tranches}

Une \textbf{tranche} permet d'extraire plusieurs éléments d'une liste.

\begin{block}{Syntaxe}
\begin{itemize}
    \item \texttt{L[m:k]} : éléments d'indice \texttt{m} à \texttt{k-1}.
    \item \texttt{L[m:]} : de l'indice \texttt{m} jusqu'à la fin.
    \item \texttt{L[:k]} : du début jusqu'à l'indice \texttt{k-1}.
\end{itemize}
\end{block}

\begin{exampleblock}{Exemple}
\begin{lstlisting}[language=Python]
nombres = [2, 3, 5, 7, 11, 13, 17, 19, 21, 23]

print(nombres[2:6])   # [5, 7, 11, 13]
print(nombres[-8:-4]) # [5, 7, 11, 13]
print(nombres[3:])    # [7, 11, 13, 17, 19, 21, 23]
print(nombres[:5])    # [2, 3, 5, 7, 11]
print(nombres[:-3])   # [2, 3, 5, 7, 11, 13, 17]
\end{lstlisting}
\end{exampleblock}

La borne de fin est \textbf{exclue} : dans \texttt{L[m:k]},
l'élément d'indice \texttt{k} n'est pas inclus.

\end{frame}


\begin{frame}[fragile]{Modifier les éléments d'une liste}

Les listes sont des objets \textbf{mutables} :
on peut modifier leur contenu après leur création.
Pour remplacer un élément, on utilise son indice :
\texttt{liste[i] = nouvelle\_valeur}.

\begin{exampleblock}{Exemple}
\begin{lstlisting}[language=Python]
capitales = ["Paris", "Berlin", "Madrid", "Athènes", "Zagreb"]
capitales[2] = "Dublin"
print(capitales)
# ['Paris', 'Berlin', 'Dublin', 'Athènes', 'Zagreb']
\end{lstlisting}
\end{exampleblock}

\vspace{0.2cm}

On peut aussi remplacer plusieurs éléments avec une tranche.

\begin{exampleblock}{Remplacer une tranche}
\begin{lstlisting}[language=Python]
capitales[2:5] = ["Dublin", "Budapest", "Vienne"]
print(capitales)
# ['Paris', 'Berlin', 'Dublin', 'Budapest', 'Vienne']
\end{lstlisting}
\end{exampleblock}

\end{frame}



\begin{frame}[fragile]{Listes en compréhension : pourquoi ?}

\textbf{Objectif :} créer une liste contenant les carrés
des entiers de $0$ à $9$.

\begin{exampleblock}{Méthode classique avec une boucle}
\begin{lstlisting}[language=Python]
carrees = []
for i in range(10):
    carrees.append(i**2)
print(carrees)
\end{lstlisting}
\end{exampleblock}

\pause

\begin{exampleblock}{Même résultat avec une compréhension}
\begin{lstlisting}[language=Python]
carrees = [i**2 for i in range(10)]
print(carrees)
\end{lstlisting}
\end{exampleblock}

\begin{block}{À retenir}
Une compréhension permet de créer une liste de façon
\textbf{plus concise}, à partir d'une séquence de valeurs.
\end{block}

\end{frame}


\begin{frame}[fragile]{Comprendre la syntaxe}

\begin{alertblock}{Syntaxe générale}
\[
\texttt{[expression for variable in sequence]}
\]
\end{alertblock}

\begin{itemize}
    \item \texttt{sequence} : les valeurs parcourues, par exemple \texttt{range(10)}.
    \item \texttt{variable} : prend successivement chaque valeur.
    \item \texttt{expression} : calcule la valeur ajoutée à la nouvelle liste.
\end{itemize}

\begin{exampleblock}{Exemple : les carrés des nombres de 0 à 9}
\begin{lstlisting}[language=Python]
carres = [x**2 for x in range(10)]
print(carres)  # [0, 1, 4, 9, 16, 25, 36, 49, 64, 81]
\end{lstlisting}
\end{exampleblock}

\pause

\begin{block}{Exercice}
Créer, avec une seule instruction, une liste contenant
la table de multiplication de 5, de $5 \times 1$ à $5 \times 10$.
\end{block}

\end{frame}
\begin{frame}[fragile]{Tester l'appartenance à une liste}

L'opérateur \texttt{in} permet de vérifier si un élément est présent
dans une liste. Le résultat est un booléen : \texttt{True} ou \texttt{False}.

\begin{exampleblock}{Exemple}
\begin{lstlisting}[language=Python]
nombres = [1, 3, 5, 7, 9, 11]
print(3 in nombres)  # True
print(2 in nombres)  # False
\end{lstlisting}
\end{exampleblock}

\pause

\begin{block}{Combiner plusieurs tests}
On peut utiliser \texttt{and/or} pour vérifier plusieurs conditions,
et \texttt{not in} pour tester l'absence d'un élément.
\begin{lstlisting}[language=Python]
print(5 in nombres and 7 in nombres)  # True
print(4 not in nombres)               # True
print(5 in nombres or 4 not in nombres)  # True
\end{lstlisting}
\end{block}

\pause

\begin{itemize}
    \item \texttt{x in liste} : vrai si \texttt{x} est présent.
    \item \texttt{x not in liste} : vrai si \texttt{x} est absent.
\end{itemize}

\end{frame}
\begin{frame}[fragile]{Copier une liste : attention à \texttt{=}}

\begin{columns}[T,onlytextwidth]
\column{0.43\textwidth}

\begin{exampleblock}{Qu'affiche ce programme?}
\begin{lstlisting}[language=Python]
L1 = [1, 2, 3, 4]
L2 = L1
L2.append(5)
print(L1)
print(L2)
\end{lstlisting}
\end{exampleblock}

\pause

\begin{alertblock}{Résultat}
\begin{lstlisting}[language=Python]
[1, 2, 3, 4, 5]
[1, 2, 3, 4, 5]
\end{lstlisting}
\end{alertblock}

\column{0.53\textwidth}
\centering
\begin{tikzpicture}[
    font=\sffamily\small,
    var/.style={draw, rounded corners, fill=gray!10,
                minimum width=1cm, minimum height=0.55cm},
    obj/.style={draw, rounded corners, fill=yellow!15,
                minimum height=0.65cm, inner sep=5pt},
    arr/.style={->, thick, >=stealth, blue!70}
]
\node[var] (v1) at (0,1.5) {\texttt{L1}};
\node[var] (v2) at (0,0.5) {\texttt{L2}};

\only<1>{
    \node[obj] (liste) at (2.3,1) {\texttt{[1, 2, 3, 4]}};
}
\only<2->{
    \node[obj] (liste) at (2.3,1) {\texttt{[1, 2, 3, 4, 5]}};
}
\draw[arr] (v1.east) -- (liste.west);
\draw[arr] (v2.east) -- (liste.west);

\only<2->{
    \node[align=center] at (1,-0.5)
    {\textcolor{red}{Même liste en mémoire}};
}
\end{tikzpicture}
\end{columns}

\vspace{0.2cm}

\begin{block}{L'idée essentielle}
\texttt{L2 = L1} ne crée pas une nouvelle liste :
les deux variables référencent  \textbf{le même objet}.
Modifier la liste via l'une modifie donc ce que l'on observe via l'autre.
\end{block}

\end{frame}


\begin{frame}[fragile]{Créer une copie avec une tranche}

Pour obtenir deux listes indépendantes, on peut utiliser
\texttt{L1[:]} : cette tranche reprend tous les éléments de \texttt{L1}.

\begin{columns}[T,onlytextwidth]
\column{0.43\textwidth}

\begin{exampleblock}{Créer puis modifier la copie}
\begin{lstlisting}[language=Python]
L1 = [1, 2, 3, 4]
L2 = L1[:]
L2.append(5)
print(L1)
print(L2)
\end{lstlisting}
\end{exampleblock}

\pause

\begin{alertblock}{Résultat}
\begin{lstlisting}[language=Python]
[1, 2, 3, 4]
[1, 2, 3, 4, 5]
\end{lstlisting}
\end{alertblock}

\column{0.53\textwidth}
\centering
\begin{tikzpicture}[
    font=\sffamily\small,
    var/.style={draw, rounded corners, fill=gray!10,
                minimum width=1cm, minimum height=0.55cm},
    obj/.style={draw, rounded corners, fill=yellow!15,
                minimum height=0.65cm, inner sep=5pt},
    arr/.style={->, thick, >=stealth, blue!70}
]
\node[var] (v1) at (0,1.7) {\texttt{L1}};
\node[var] (v2) at (0,0.3) {\texttt{L2}};

\node[obj] (liste1) at (2.5,1.7) {\texttt{[1, 2, 3, 4]}};

\only<1>{
    \node[obj] (liste2) at (2.5,0.3) {\texttt{[1, 2, 3, 4]}};
}
\only<2->{
    \node[obj] (liste2) at (2.5,0.3) {\texttt{[1, 2, 3, 4, 5]}};
}

\draw[arr] (v1.east) -- (liste1.west);
\draw[arr] (v2.east) -- (liste2.west);

\only<2->{
    \node[align=center] at (1,-0.55)
    {\textcolor{blue!70}{Deux listes distinctes}};
}
\end{tikzpicture}
\end{columns}

\vspace{0.2cm}

\texttt{L2 = L1[:]} crée une nouvelle liste contenant les mêmes éléments.
Ajouter ou supprimer un élément de \texttt{L2} ne modifie pas \texttt{L1}.

\end{frame}


\begin{frame}[fragile]{Une autre méthode : \texttt{list()}}

La fonction \texttt{list()} permet également de créer une copie
superficielle d'une liste.

\begin{columns}[T,onlytextwidth]
\column{0.43\textwidth}

\begin{exampleblock}{Exemple}
\begin{lstlisting}[language=Python]
L1 = [1, 2, 3, 4]
L2 = list(L1)
L2.append(5)
print(L1)
print(L2)
\end{lstlisting}
\end{exampleblock}

\begin{alertblock}{Résultat}
\begin{lstlisting}[language=Python]
[1, 2, 3, 4]
[1, 2, 3, 4, 5]
\end{lstlisting}
\end{alertblock}

\column{0.53\textwidth}
\centering
\begin{tikzpicture}[
    font=\sffamily\small,
    var/.style={draw, rounded corners, fill=gray!10,
                minimum width=1cm, minimum height=0.55cm},
    obj/.style={draw, rounded corners, fill=yellow!15,
                minimum height=0.65cm, inner sep=5pt},
    arr/.style={->, thick, >=stealth, blue!70}
]
\node[var] (v1) at (0,1.7) {\texttt{L1}};
\node[var] (v2) at (0,0.3) {\texttt{L2}};

\node[obj] (liste1) at (2.5,1.7) {\texttt{[1, 2, 3, 4]}};
\node[obj] (liste2) at (2.5,0.3) {\texttt{[1, 2, 3, 4, 5]}};

\draw[arr] (v1.east) -- (liste1.west);
\draw[arr] (v2.east) -- (liste2.west);
\end{tikzpicture}
\end{columns}

\vspace{0.2cm}

\begin{block}{Deux façons de copier}
\centering
\texttt{L2 = L1[:]} \qquad ou \qquad \texttt{L2 = list(L1)}
\end{block}

\end{frame}
\begin{frame}[fragile]{Concaténer deux listes}

L'opérateur \texttt{+} permet de \textbf{rassembler deux listes}
pour former une nouvelle liste.

\begin{exampleblock}{Que va afficher ce programme ?}
\begin{lstlisting}[language=Python]
capitales1 = ["Paris", "Berlin", "Madrid"]
capitales2 = ["Bruxelles", "Oslo"]

capitales1 + capitales2
print(capitales1)
\end{lstlisting}
\end{exampleblock}

\pause

\begin{alertblock}{Attention !}
L'expression \texttt{capitales1 + capitales2} construit une nouvelle liste,
mais ne modifie pas \texttt{capitales1} si le résultat n'est pas conservé.
\end{alertblock}

\pause

\begin{exampleblock}{Créer une nouvelle liste}
\begin{lstlisting}[language=Python]
capitales3 = capitales1 + capitales2
print(capitales3)
# ['Paris', 'Berlin', 'Madrid', 'Bruxelles', 'Oslo']
\end{lstlisting}
\end{exampleblock}

\end{frame}


\begin{frame}[fragile]{Ajouter les éléments d'une liste à une autre}

Si l'objectif est de modifier une liste existante, on peut utiliser
\texttt{extend()}.

\begin{exampleblock}{Exemple}
\begin{lstlisting}[language=Python]
capitales1 = ["Paris", "Berlin", "Madrid"]
capitales2 = ["Bruxelles", "Oslo"]

capitales1.extend(capitales2)
print(capitales1)
# ['Paris', 'Berlin', 'Madrid', 'Bruxelles', 'Oslo']
\end{lstlisting}
\end{exampleblock}

\begin{block}{À retenir}
\begin{itemize}
    \item \texttt{L3 = L1 + L2} : crée une nouvelle liste.
    \item \texttt{L1.extend(L2)} : ajoute les éléments de \texttt{L2} à \texttt{L1}.
\end{itemize}
\end{block}

\end{frame}




\begin{frame}[fragile]{Listes imbriquées : accéder aux éléments}

Une \textbf{liste imbriquée} est une liste dont certains éléments
sont eux-mêmes des listes.

\begin{exampleblock}{Exemple}
\begin{lstlisting}[language=Python]
L = [[1, 2, 3], [4, 5], [6, 7, 8, 9]]
\end{lstlisting}
\end{exampleblock}

\begin{center}
\begin{tikzpicture}[
    font=\sffamily,
    listbox/.style={draw, rounded corners, fill=blue!8,
                    minimum height=0.7cm, inner sep=5pt},
    item/.style={draw, rounded corners, fill=yellow!15,
                 minimum height=0.6cm, inner sep=5pt}
]
\node[listbox] (l0) at (0,1.5) {\texttt{L[0]}};
\node[item] (a) at (2,1.5) {\texttt{[1, 2, 3]}};

\node[listbox] (l1) at (0,0.5) {\texttt{L[1]}};
\node[item] (b) at (2,0.5) {\texttt{[4, 5]}};

\node[listbox] (l2) at (0,-0.5) {\texttt{L[2]}};
\node[item] (c) at (2,-0.5) {\texttt{[6, 7, 8, 9]}};

\draw[->, thick] (l0.east) -- (a.west);
\draw[->, thick] (l1.east) -- (b.west);
\draw[->, thick] (l2.east) -- (c.west);
\end{tikzpicture}
\end{center}

\pause

\begin{exampleblock}{Accéder à une liste ou à un élément}
\begin{lstlisting}[language=Python]
print(L[0])     # [1, 2, 3]
print(L[1])     # [4, 5]
print(L[2][1])  # 7
\end{lstlisting}
\end{exampleblock}

\begin{block}{À retenir}
\texttt{L[i]} désigne la liste d'indice \texttt{i}.
\texttt{L[i][j]} désigne l'élément d'indice \texttt{j}
dans cette liste.
\end{block}

\end{frame}


\begin{frame}[fragile]{Parcourir une liste imbriquée}

On veut calculer la somme de tous les nombres contenus dans \texttt{L}.

\begin{exampleblock}{Deux boucles imbriquées}
\begin{lstlisting}[language=Python]
L = [[1, 2, 3], [4, 5], [6, 7, 8, 9]]
somme = 0
for sous_liste in L:
    for nombre in sous_liste:
        somme += nombre
print(somme)  # 45
\end{lstlisting}
\end{exampleblock}

\pause

\begin{block}{Comment fonctionne le parcours ?}
\begin{enumerate}
    \item La première boucle parcourt chaque sous-liste.
    \item La seconde boucle parcourt chaque nombre de cette sous-liste.
    \item Chaque nombre est ajouté à \texttt{somme}.
\end{enumerate}
\end{block}

\pause

La liste contient trois sous-listes, mais elles n'ont pas toutes
la même longueur. Les boucles \texttt{for} parcourent chaque élément
sans avoir besoin de connaître ces longueurs à l'avance.
\end{frame}




\begin{frame}[fragile]{Créer une liste de listes : attention aux références}

On veut créer une liste de $m$ lignes contenant chacune $n$ zéros.

\begin{exampleblock}{Une première tentative}
\begin{lstlisting}[language=Python]
m, n = 3, 4
L = [[0] * n] * m
\end{lstlisting}
\end{exampleblock}

Cette instruction crée bien $m$ références vers des lignes, mais
\textbf{toutes désignent la même liste}.

\begin{columns}[T,onlytextwidth]
\column{0.42\textwidth}

\begin{exampleblock}{Modifier un élément}
\begin{lstlisting}[language=Python]
L[0][1] = 1
\end{lstlisting}
\end{exampleblock}

\pause

\begin{alertblock}{Résultat}
\begin{lstlisting}[language=Python]
[[0, 1, 0, 0],
 [0, 1, 0, 0],
 [0, 1, 0, 0]]
\end{lstlisting}
\end{alertblock}

\column{0.54\textwidth}
\centering

\begin{tikzpicture}[
    font=\sffamily\small,
    frame/.style={
        draw, fill=gray!10, rounded corners,
        minimum width=0.65cm, minimum height=0.55cm
    },
    cell/.style={
        draw, fill=blue!8, minimum width=0.8cm,
        minimum height=0.65cm
    },
    object/.style={
        draw, fill=yellow!15, rounded corners,
        inner sep=5pt
    },
    arr/.style={->, thick, >=stealth, blue!65}
]

% Variable L
\node[frame] (varL) at (0,2.1) {\texttt{L}};

% Les trois références dans L
\node[cell] (ref0) at (1.5,2.8) {\texttt{0}};
\node[cell] (ref1) at (1.5,2.1) {\texttt{1}};
\node[cell] (ref2) at (1.5,1.4) {\texttt{2}};

% Indices au-dessus des cases
\node[font=\scriptsize] at (1.5,3.25) {Indices};

% La même sous-liste est partagée

\node[object] (row) at (4,2.1) {
    \texttt{[0, \textcolor{red}{1}, 0, 0]}
};

% L pointe vers la liste contenant les références
\draw[arr] (varL.east) -- (1.0,2.1);

% Chaque case pointe vers la même sous-liste
\draw[arr] (ref0.east) to[out=0,in=150] (row.north west);
\draw[arr] (ref1.east) -- (row.west);
\draw[arr] (ref2.east) to[out=0,in=210] (row.south west);

% Explication
\node[align=center, text width=4.7cm] at (2,0.45)
    {Les trois cases référencent la même sous-liste.};

\end{tikzpicture}
\end{columns}
\end{frame}


\begin{frame}[fragile]{La bonne méthode : créer chaque ligne séparément}

Une compréhension de liste permet de créer une nouvelle ligne à chaque
itération.

\begin{exampleblock}{Créer une grille de zéros}
\begin{lstlisting}[language=Python]
m, n = 3, 4
L = [[0] * n for i in range(m)]
\end{lstlisting}
\end{exampleblock}

\begin{columns}[T,onlytextwidth]
\column{0.43\textwidth}
\begin{exampleblock}{Modifier un élément}
\begin{lstlisting}[language=Python]
L[0][1] = 1
print(L)
\end{lstlisting}
\end{exampleblock}

\pause

\begin{alertblock}{Résultat}
\begin{lstlisting}[language=Python]
[[0, 1, 0, 0],
 [0, 0, 0, 0],
 [0, 0, 0, 0]]
\end{lstlisting}
\end{alertblock}

\column{0.54\textwidth}
\centering
\begin{tikzpicture}[
    font=\sffamily\small,
    frame/.style={
        draw, fill=gray!10, rounded corners,
        minimum width=0.7cm, minimum height=0.55cm
    },
    cell/.style={
        draw, fill=blue!8,
        minimum width=0.7cm, minimum height=0.65cm
    },
    object/.style={
        draw, fill=yellow!15, rounded corners,
        inner sep=5pt
    },
    arr/.style={->, thick, >=stealth, blue!65}
]

% Variable L
\node[frame] (varL) at (0,1.5) {\texttt{L}};

% Liste contenant les références vers les lignes
\node[cell] (ref0) at (1.5,2.2) {\texttt{0}};
\node[cell] (ref1) at (1.5,1.5) {\texttt{1}};
\node[cell] (ref2) at (1.5,0.8) {\texttt{2}};

% Les trois sous-listes indépendantes
\node[object] (r0) at (4,2.2) {\texttt{[0, 0, 0, 0]}};
\node[object] (r1) at (4,1.5) {\texttt{[0, 0, 0, 0]}};
\node[object] (r2) at (4,0.8) {\texttt{[0, 0, 0, 0]}};

% L désigne la liste de références
\draw[arr] (varL.east) -- (ref1.west);

% Chaque case pointe vers sa propre sous-liste
\draw[arr] (ref0.east) -- (r0.west);
\draw[arr] (ref1.east) -- (r1.west);
\draw[arr] (ref2.east) -- (r2.west);

% Explication
\node[align=center, text width=5cm] at (2.8,0)
    {Chaque position désigne une ligne indépendante.};

\end{tikzpicture}
\end{columns}

\begin{block}{À retenir}
\texttt{[[0] * n for i in range(m)]} crée $m$ listes distinctes,
chacune contenant $n$ zéros.
\end{block}

\end{frame}


\begin{frame}[fragile]{Exercice : construire une matrice}

Écrire une fonction \texttt{construire\_matrice(n)} qui construit
une matrice carrée de taille $n$ suivant la règle :

\[
L_{i,j} =
\begin{cases}
1 & \text{si } i=j,\\
2 & \text{si } i<j,\\
0 & \text{si } i>j.
\end{cases}
\]

Pour $n=4$, on veut obtenir :

\[
\begin{bmatrix}
1 & 2 & 2 & 2\\
0 & 1 & 2 & 2\\
0 & 0 & 1 & 2\\
0 & 0 & 0 & 1
\end{bmatrix}
\]

\begin{block}{Consignes}
\begin{itemize}
    \item Créer d'abord une matrice remplie de zéros.
    \item Compléter les éléments en parcourant les lignes et les colonnes.
    \item La fonction doit \textbf{renvoyer} la matrice avec \texttt{return}.
\end{itemize}
\end{block}

\end{frame}


\begin{frame}[fragile]{Correction : construire une matrice}

\begin{exampleblock}{Une solution avec deux boucles}
\begin{lstlisting}[language=Python]
def construire_matrice(n):
    L = [[0] * n for i in range(n)]
    for i in range(n):
        for j in range(n):
            if i == j:
                L[i][j] = 1
            elif i < j:
                L[i][j] = 2
    return L
print(construire_matrice(4))
\end{lstlisting}
\end{exampleblock}

\pause

\begin{block}{Pourquoi cette solution fonctionne-t-elle ?}
\begin{itemize}
    \item La compréhension crée $n$ lignes indépendantes.
    \item Les deux boucles parcourent chaque position $(i,j)$.
    \item Les conditions déterminent la valeur à placer.
    \item \texttt{return L} renvoie la matrice construite.
\end{itemize}
\end{block}

\end{frame}


\begin{frame}[fragile]{Une autre solution : construire chaque ligne}

On peut aussi construire chaque ligne directement, sans modifier
chaque élément séparément.

\begin{exampleblock}{Correction compacte}
\begin{lstlisting}[language=Python]
def construire_matrice(n):
    L = []
    for i in range(n):
        ligne = [0] * i + [1] + [2] * (n - i - 1)
        L.append(ligne)
    return L
print(construire_matrice(4))
\end{lstlisting}
\end{exampleblock}

\pause

\begin{block}{Comprendre la construction d'une ligne}
Pour la ligne d'indice \texttt{i}, on assemble :
\begin{itemize}
    \item \texttt{i} zéros ;
    \item un \texttt{1} sur la diagonale ;
    \item \texttt{n - i - 1} deux après la diagonale.
\end{itemize}
\end{block}

\end{frame}








\begin{comment}
\begin{frame}[fragile]{Les listes}

Un \textcolor{blue}{tableau} est une structure de données contenant une série
d'éléments, où la position de chaque élément compte.\\

%\vspace{0.5 cm}
\begin{alertblock}{Syntaxe}

On peut accéder aux éléments d'une liste par leur \textbf{indice}.
\begin{verbatim}
tableau  : ["bus", "voiture", "vélo", "trotinette"]
indice :    0        1        2         3

\end{verbatim}
\end{alertblock}

\pause

\vspace{0.3 cm}
En Python, on utilise le type \texttt{\texttt{list}} pour les représenter. 
%En Python, une liste est créée en plaçant des éléments entre \textbf{crochets} \texttt{{[}{]}} et en les séparant par des \emph{virgules}.

\begin{exampleblock}{}
\small{
\begin{verbatim}
# Exemple: Liste des capitales européennes
capitales = ["Paris", "Berlin", "Madrid", "Athènes", "Zagreb", "Rome"]

# Exemple: Liste d'entiers
entiers = [5, 0, 17, 42, 23]
\end{verbatim}
}
\end{exampleblock}
\end{frame}

\frame[containsverbatim]{
\frametitle{Les listes}

 La construction de listes en Python est très \emph{flexible}.
\begin{itemize}
\item Il est possible de créer des listes contenant des valeurs de
\textbf{types différents}.
\end{itemize}

\begin{exampleblock}{}
\begin{verbatim}

# Liste mixte
mixte = [7, "tennis", True, 5.3]

\end{verbatim}
\end{exampleblock}
\small{Cette liste contient à la fois des valeurs de type \texttt{int},
\texttt{string}, \texttt{boolean} et \texttt{float}.}


\vspace{0.5 cm}


\begin{itemize}
\item La fonction \texttt{len()} nous permet de connaître le nombre
  d'éléments dans une liste.
\end{itemize}

\begin{exampleblock}{}
\begin{verbatim}
liste = [0, 1, 2, 3, 4, 5, 6, 7, 8, 9]
print("La liste contient", len(liste), "éléments.")
\end{verbatim}
\end{exampleblock}





}

\frame[containsverbatim]{
\frametitle{Longueur d'une liste}




\begin{itemize}
\item Une liste a une longueur variable. La méthode \texttt{append()} permet d'ajouter un élément à la fin d'une
liste.
\end{itemize}


\begin{exampleblock}{}
\begin{verbatim}
mixte = [7, "tennis", True, 5.3]

mixte.append("natation")
print(mixte)

\end{verbatim}
\end{exampleblock}

\begin{itemize}
\item  La méthode \texttt{remove()} permet d'enlever un élément  de la liste. 
\end{itemize}

\begin{exampleblock}{}
\begin{verbatim}

mixte.remove("tennis")
mixte.remove(7)
print(mixte)

\end{verbatim}
\end{exampleblock}

\small{Attention: les méthodes \texttt{append} et \texttt{remove} sont propres à l'objet (\texttt{mixte}). Chose à clarifier plus tard ...}

}




\frame[containsverbatim]{
\frametitle{Accéder aux éléments d'une liste}


 Pour une liste de longueur \(n\)   :


\begin{itemize}

\item   L'indice est la \textbf{position} d'un élément dans la liste.
\item  Pour une liste de \(n\) éléments, les indices vont de 0 à \(n-1\).
\item Les indices doivent être des nombres \textbf{entiers}.

 \end{itemize}

\begin{exampleblock}{Un premier exemple}
\begin{verbatim}
liste = ["bus", "voiture", "velo", "trotinette"]

print(liste[0])
print(liste[2])
\end{verbatim}
\end{exampleblock}


}

\frame[containsverbatim]{
\frametitle{Indiçage négatif}

 \small{Chaque élément d'une liste  a deux indices :
   un indice positif \(i\) et un indice négatif \(i - n\).}

\begin{alertblock}{}
\begin{verbatim}
liste          : ["bus", "voiture", "velo", "trotinette"]
indice positif :    0       1         2         3
indice négatif :   -4      -3        -2        -1 
\end{verbatim}
\end{alertblock}

\begin{exampleblock}{Indiçage négatif}
\begin{verbatim}
liste = ["bus", "voiture", "vélo", "trotinette"]
print(liste[-1])
print(liste[-2])
print(liste[0],liste[-4])
\end{verbatim}
\end{exampleblock}

\vspace{0.5 cm}

    \begin{itemize}

\item
  L'indiçage négatif revient à compter de la fin.
\item
  \textbf{Avantage} : Accèder au dernier ou avant-dernier élément de la
  liste, sans connaître sa longueur (indices -1 et -2 respectivement)
\end{itemize}

}


\frame[containsverbatim]{
\frametitle{Tranches}

On peut accéder à une partie (une tranche) d'une liste L de longueur $n$ en utilisant des indices :

\begin{itemize}
\item \texttt{L[m:k]} donne accès aux éléments \texttt{L[m], L[m+1], ... ,L[k-1]}
\item \texttt{L[m:]} donne accès aux éléments \texttt{L[m], ..., L[n-1]}
\item  \texttt{L[:k]} donne accès aux éléments \texttt{L[0], L[1], ... ,L[k-1]}  
\end{itemize}

\vspace{0.5 cm}
\begin{exampleblock}{Exemple}
\begin{verbatim}
liste = [2, 3, 5, 7, 11, 13, 17, 19, 21, 23]
print(liste[2:6])
print(liste[-8:-4])
liste[3:]
liste[:5]
liste[:-3]
\end{verbatim}
\end{exampleblock}

}


\frame[containsverbatim]{
\frametitle{Modifier ou ajouter des éléments dans une
liste}


\textbf{Rappel.} Les listes sont des objets \textcolor{blue}{mutables}. Leurs éléments peuvent
donc être modifiés.


\begin{itemize}
\item En écrivant \texttt{liste{[}i{]}\ =\ nouvel\_element}, l'élément de la
liste à l'indice \texttt{i} est replacé par \texttt{nouvel\_element}.
 \end{itemize}

\vspace{1 cm}

\begin{exampleblock}{}
\begin{verbatim}
capitales = ["Paris", "Berlin", "Madrid", "Athènes", "Zagreb"]
capitales[2:5] = ["Dublin", "Budapest", "Vienne"]
print(capitales)
\end{verbatim}
\end{exampleblock}





}




\frame[containsverbatim]{
\frametitle{Listes en compréhension}



\vspace{0.5 cm}

Imaginons, qu'on veut créer une liste contenant la racine carrée des
entiers \(0, 1, \dots, 9\).

\vspace{0.5 cm}

Avec une boucle: 
\begin{exampleblock}{}
\begin{verbatim}
import math
l = []
for i in range(10) :
    l.append(math.sqrt(i))

print(l)
\end{verbatim}
\end{exampleblock}

\vspace{0.5 cm}

De manière compacte, en une seule instruction:
\begin{exampleblock}{}
\begin{verbatim}
l = [math.sqrt(i) for i in range(10)]

print(l)

\end{verbatim}
\end{exampleblock}

}



\begin{frame}[fragile]{Listes en compréhension}

On peut créer des listes de façon \textbf{simple}, \textbf{concise} et
\textbf{compacte}.\\


\begin{alertblock}{Syntaxe}
%Une liste \texttt{L} en compréhension se créé de la façon suivante :
\texttt{L} = \texttt{{[}expr\ for\ x\ in\ t{]}}
\end{alertblock}

\begin{itemize}
\item
  Dans cette syntaxe : les crochets \texttt{{[}{]}}, les mots clés
  \texttt{for} et \texttt{in} sont obligatoires.
\item
  \texttt{t} est souvent une liste ou une collection de type
  \texttt{range}.
\item
  \texttt{x} est la variable de contrôle
\item
  \texttt{expr} est une expression qui dépend en général de \texttt{x}
  et dont la valeur est placée dans la liste.
\end{itemize}


\vspace{0.5 cm}

Exercice : Créer avec une seule instruction une liste qui contient la table de multiplication de 5.

\pause
\begin{exampleblock}{}
\begin{verbatim}
mult5 = [5*i for i in range(1,11)]
\end{verbatim}
\end{exampleblock}
\end{frame}

\begin{frame}[fragile]{Test d’appartenance}

On peut facilement tester l'appartenance d'un élément dans une liste avec l'opérateur \texttt{in}.

\begin{exampleblock}{}
\begin{verbatim}
nombres = [1, 3, 5, 7, 9, 11]
print(3 in nombres)
print(2 in nombres)
print(5 in nombres and 7 in nombres)
\end{verbatim}
\end{exampleblock}

\vspace{1 cm}

En combinant avec le mot-clé \texttt{not} on peut vérifier si un élément est absent d'une liste.

\begin{exampleblock}{}
\begin{verbatim}
4 not in nombres
\end{verbatim}
\end{exampleblock}

\end{frame}



\begin{frame}[fragile]{Copie d'une liste}

Qu'affiche-t-on à l'exécution? 

\begin{columns}
\column{4 cm}
\begin{exampleblock}{}
\begin{verbatim}
L1 = [1, 2, 3, 4]
L2 = L1
print(L2)

L2.append(5)
print(L2)
print(L1) 
\end{verbatim}
\end{exampleblock}
\column{5 cm}
\begin{tikzpicture}[
    font=\sffamily,
    frame/.style={draw, fill=gray!10, rounded corners, inner sep=1pt},
    object/.style={draw, fill=yellow!10, rounded corners, inner sep=0.5pt},
    arr/.style={->, thick, >=stealth, blue!60},
    node distance=0.3cm and 0.5cm
]

% Frames box
\node[frame, align=left, minimum width=0.3cm] (frame) at (3,0.5){
       \begin{tabular}{l }
        L1  \\
    \end{tabular}
};
\node[frame, align=left, minimum width=0.3cm] (frame1) at (3,1.5){
       \begin{tabular}{l }
        L2  \\
    \end{tabular}
};


\node[object, anchor=west] (Llist)  at (4,3.1) {
    \begin{tabular}{ccccc}        
        1 & 2 & 3 & 4 & \\
    \end{tabular}
};



% Arrows from L list entries to the row lists
%\draw[arr] ([yshift=1mm]Llist.north west) ++(0.3cm,-0.1) .. controls +(-0.5,1) and +(-0.5,-0.5) .. (list2.west);

% Arrow from L in frame to Llist
\draw[arr] ([xshift=0.6cm]frame.west)++(0.4cm,0) .. controls +(0.5,0.5) and +(-0.4,-0.4) .. (Llist.west);

\draw[arr] ([xshift=0.6cm]frame1.west)++(0.37cm,0) .. controls +(0.7,0.5) and +(-1,0) .. (Llist.west);

\only<2>
\node[object, anchor=west] (list3)  at (4,3.1) {
    \begin{tabular}{ccccccc}        
        1 & 2 & 3 & 4 & \textcolor{red}{5} \\
    \end{tabular}
};


\end{tikzpicture}



\end{columns}

}

\vspace{0.5 cm}

\begin{itemize}
\item
  Il ne faut jamais utiliser l'opérateur \texttt{=} pour copier une
  liste.
\item
  Les deux objets \texttt{L1} et \texttt{L2} partagent la \textbf{même
  zone mémoire}. Une modification de l'un entraine une modification de
  l'autre.
\end{itemize}

\end{frame}

\begin{frame}[fragile]{Copie d'une liste}


\begin{columns}
\column{4 cm}
\begin{exampleblock}{}
\begin{verbatim}
L1 = [1, 2, 3, 4]
L2 = L1[:]
print(L2)

L2.append(5)
print(L2)
print(L1)
\end{verbatim}
\end{exampleblock}
\column{5 cm}
\begin{tikzpicture}[
    font=\sffamily,
    frame/.style={draw, fill=gray!10, rounded corners, inner sep=1pt},
    object/.style={draw, fill=yellow!10, rounded corners, inner sep=0.5pt},
    arr/.style={->, thick, >=stealth, blue!60},
    node distance=0.3cm and 0.5cm
]

% Frames box
\node[frame, align=left, minimum width=0.3cm] (frame) at (3,0.5){
       \begin{tabular}{l }
        L2  \\
    \end{tabular}
};
\node[frame, align=left, minimum width=0.3cm] (frame1) at (3,1.5){
       \begin{tabular}{l }
        L1  \\
    \end{tabular}
};


\node[object, anchor=west] (Llist)  at (4.5, 1) {
    \begin{tabular}{cccc}        
        1 & 2 & 3 & 4  \\
    \end{tabular}
};


\node[object, anchor=west] (Llist1)  at (4.5,3.1.5) {
    \begin{tabular}{ccccc}        
        1 & 2 & 3 & 4 & \textcolor{red}{5}\\
    \end{tabular}
};


% Arrows from L list entries to the row lists
%\draw[arr] ([yshift=1mm]Llist.north west) ++(0.3cm,-0.1) .. controls +(-0.5,1) and +(-0.5,-0.5) .. (list2.west);

% Arrow from L in frame to Llist
\draw[arr] ([xshift=0.6cm]frame.west)++(0.4cm,0) .. controls +(0.5,0.5) and +(-0.4,-0.4) .. (Llist1.west);

\draw[arr] ([xshift=0.6cm]frame1.west)++(0.37cm,0) .. controls +(0.7,0.5) and +(-1,0) .. (Llist.west);


\end{tikzpicture}



\end{columns}
\end{frame}

\begin{frame}[fragile]{Copie d'une liste}

On peut utiliser la fonction $\texttt{list()}$ qui permet de créer une nouvelle liste. 


\begin{columns}
\column{5 cm}
\begin{exampleblock}{}
\begin{verbatim}
L1 = [1, 2, 3, 4]
L2 = list(L1)
print(L2)

L2.append(5)
print(L2)
print(L1)

\end{verbatim}
\end{exampleblock}
\column{5 cm}
\begin{tikzpicture}[
    font=\sffamily,
    frame/.style={draw, fill=gray!10, rounded corners, inner sep=1pt},
    object/.style={draw, fill=yellow!10, rounded corners, inner sep=0.5pt},
    arr/.style={->, thick, >=stealth, blue!60},
    node distance=0.3cm and 0.5cm
]

% Frames box
\node[frame, align=left, minimum width=0.3cm] (frame) at (3,0.5){
       \begin{tabular}{l }
        L2  \\
    \end{tabular}
};
\node[frame, align=left, minimum width=0.3cm] (frame1) at (3,1.5){
       \begin{tabular}{l }
        L1  \\
    \end{tabular}
};


\node[object, anchor=west] (Llist)  at (4.5, 1) {
    \begin{tabular}{cccc}        
        1 & 2 & 3 & 4  \\
    \end{tabular}
};


\node[object, anchor=west] (Llist1)  at (4.5,3.1.5) {
    \begin{tabular}{ccccc}        
        1 & 2 & 3 & 4 & \textcolor{red}{5}\\
    \end{tabular}
};


% Arrows from L list entries to the row lists
%\draw[arr] ([yshift=1mm]Llist.north west) ++(0.3cm,-0.1) .. controls +(-0.5,1) and +(-0.5,-0.5) .. (list2.west);

% Arrow from L in frame to Llist
\draw[arr] ([xshift=0.6cm]frame.west)++(0.4cm,0) .. controls +(0.5,0.5) and +(-0.4,-0.4) .. (Llist1.west);

\draw[arr] ([xshift=0.6cm]frame1.west)++(0.37cm,0) .. controls +(0.7,0.5) and +(-1,0) .. (Llist.west);


\end{tikzpicture}
\end{columns}

\end{frame}

\begin{frame}[fragile]{Concatener deux listes}

Qu'affiche-t-on à l'exécution? 

\vspace{0.5 cm}

\begin{exampleblock}{Exemple (rappel)}
\small{
\begin{verbatim}
capitales1 = ["Paris", "Berlin", "Madrid", "Athènes", "Zagreb", "Rome"]
capitales2 = ["Bruxelles", "Oslo"]

capitale1=capitales1 + capitales2
print(capitales1)  
\end{verbatim}
}
\end{exampleblock}

\vspace{0.5 cm}

Comment faire pour créer une nouvelle liste contenant la concatenation? 

\end{frame}


\begin{frame}[fragile]{Listes imbriquées}

Une liste où chaque élément est à nouveau une liste s'appelle 
\textbf{liste imbriquée}, \textbf{double liste} ou
\textbf{liste bidimensionnelle}.


\vspace{0.5 cm}

\begin{exampleblock}{}
\begin{verbatim}
L = [[1, 2, 3], [4, 5], [6, 7, 8, 9]]
print(L[0])
print(L[1])
print(L[2][1])
\end{verbatim}
\end{exampleblock}
\end{frame}


\begin{frame}[fragile]{Listes imbriquées}

On veut  calculer la somme de tous les éléments de la liste
\texttt{L}.


\begin{exampleblock}{}
\begin{verbatim}
L = [[1, 2, 3], [4, 5], [6, 7, 8, 9]]

somme = 0
for i in range(len(L)) :
    for j in range(len(L[i])) :
        somme += L[i][j]

print("La somme de tous les éléments de la liste vaut :", somme)
\end{verbatim}
\end{exampleblock}

\pause 

ou encore 


\begin{exampleblock}{}
\begin{verbatim}
somme = 0
for i in L :
    for j in i :
        somme += j
        
print("La somme de tous les éléments de la liste vaut :", somme)  
\end{verbatim}
\end{exampleblock}



\end{frame}

\begin{frame}[fragile]{Création de listes imbriquées}


\begin{exampleblock}{}
\begin{verbatim}
m = 3
n = 4
L = [[0] * n] * m
print(L)
\end{verbatim}
\end{exampleblock}
\texttt{m} listes sont créées mais qui font toutes référence à la même liste !

\vspace{0.5 cm}

\begin{columns}
\column{3 cm}
\begin{exampleblock}{}
\begin{verbatim}
L[0][1] = 1
print(L)
\end{verbatim}
\end{exampleblock}
%\vspace{0.5 cm}
%Comment faire de manière en évitant une copie de références ?
\column{5 cm}
\begin{tikzpicture}[
    font=\sffamily,
    frame/.style={draw, fill=gray!10, rounded corners, inner sep=1pt},
    object/.style={draw, fill=yellow!10, rounded corners, inner sep=0.5pt},
    arr/.style={->, thick, >=stealth, blue!60},
    node distance=0.3cm and 0.5cm
]

% Frames box
\node[frame, align=left, minimum width=0.3cm] (frame) at (3,0.5){
       \begin{tabular}{l }
        L  \\
    \end{tabular}
};
% List of lists (L)
\node[object, minimum width=0.5cm] (Llist)  at (4,2) {
    \begin{tabular}{c  |  c | c}
    $~ ^{\tiny{\texttt{0}}}$ & $~ ^{\tiny{\texttt{1}}}$ & $~ ^{\tiny{\texttt{2}}}$ \\     
    \end{tabular}
};

\node[object, anchor=west] (list2)  at (5,3.7) {
    \begin{tabular}{cccc}        
        0 & 0 & 0 & 0
    \end{tabular}
};


% Arrows from L list entries to the row lists
\draw[arr] ([yshift=1mm]Llist.north west) ++(0.3cm,-0.1) .. controls +(-0.5,1) and +(-0.5,-0.5) .. (list2.west);
\draw[arr] ([yshift=1mm]Llist.north west) ++(1cm,-0.1) .. controls +(-0.3,0.8) and +(-0.5,-0.5) .. (list2.west);
\draw[arr] ([yshift=1mm]Llist.north west) ++(1.5cm,-0.1) .. controls +(0,0.5) and +(-0.5,-0.5) .. (list2.west);

% Arrow from L in frame to Llist
\draw[arr] ([xshift=0.6cm]frame.west)++(0.25cm,0) .. controls +(0.5,0.5) and +(-1,0) .. (Llist.west);

\only<2>
\node[object, anchor=west] (list3)  at (5,3.7) {
    \begin{tabular}{cccc}        
        0 & \textcolor{red}{1} & 0 & 0
    \end{tabular}
};

\end{tikzpicture}
\end{columns}
\end{frame}

\begin{frame}[fragile]{Création de listes imbriquées}

Plusieurs façons correctes de faire. La plus compacte est sûrement la suivante :

\begin{exampleblock}{}
\begin{verbatim}
m = 3
n = 4
L = [[0] * n for i in range(m)]
print(L)
\end{verbatim}
\end{exampleblock}

\vspace{1 cm}

\begin{columns}
\column{6 cm}
\begin{tikzpicture}[
    font=\sffamily,
    frame/.style={draw, fill=gray!10, rounded corners, inner sep=1pt},
    object/.style={draw, fill=yellow!10, rounded corners, inner sep=0.5pt},
    arr/.style={->, thick, >=stealth, blue!60},
    node distance=0.3cm and 0.5cm
]

% Frames box
\node[frame, align=left, minimum width=0.3cm] (frame) at (3,0.5){
       \begin{tabular}{l }
        L  \\
    \end{tabular}
};
% List of lists (L)
\node[object, minimum width=0.5cm] (Llist)  at (4,2) {
    \begin{tabular}{c  |  c | c}
    $~ ^{\tiny{\texttt{0}}}$ & $~ ^{\tiny{\texttt{1}}}$ & $~ ^{\tiny{\texttt{2}}}$ \\     
    \end{tabular}
};

% Three 1D lists (rows)
\node[object, anchor=west] (list1) at (5,3) {
    \begin{tabular}{cccc}
             0 & 0 & 0 & 0
    \end{tabular}
};

\node[object, anchor=west] (list2)  at (5,3.7) {
    \begin{tabular}{cccc}        
        0 & 0 & 0 & 0
    \end{tabular}
};

\node[object, anchor=west] (list3)  at (5,4.4){
    \begin{tabular}{cccc}    
        0 & 0 & 0 & 0
    \end{tabular}
};

% Arrows from L list entries to the row lists
\draw[arr] ([yshift=1mm]Llist.north west) ++(0.3cm,-0.1) .. controls +(-0.5,1) and +(-0.5,-0.5) .. (list3.west);
\draw[arr] ([yshift=1mm]Llist.north west) ++(1cm,-0.1) .. controls +(-0.3,0.8) and +(-0.5,-0.5) .. (list2.west);
\draw[arr] ([yshift=1mm]Llist.north west) ++(1.5cm,-0.1) .. controls +(0,0.5) and +(-0.5,-0.5) .. (list1.west);

% Arrow from L in frame to Llist
\draw[arr] ([xshift=0.6cm]frame.west)++(0.25cm,0) .. controls +(0.5,0.5) and +(-1,0) .. (Llist.west);

\only<2->
\node[object, anchor=west] (list4) at (5,4.4) {
    \begin{tabular}{cccc}
             0 & \textcolor{red}{1} & 0 & 0
    \end{tabular}
};


\end{tikzpicture}


\column{3 cm}
\begin{exampleblock}{}
\begin{verbatim}
L[0][1] = 1
print(L)
\end{verbatim}
\end{exampleblock}
\end{columns}

\end{frame}






\begin{frame}[fragile]{Création de listes imbriquées}

\textbf{Exercice} : Créer une liste en deux dimensions, contenant la
matrice suivante :

\[\begin{bmatrix}1 & 2 & 2 & 2 \\0 & 1 & 2 & 2 \\ 0 & 0 & 1 & 2 \\ 0 & 0 & 0 & 1\end{bmatrix}\]

\pause

\begin{exampleblock}{}
\begin{verbatim}
L = [[0] * 4 for i in range(4)]
for i in range(4) :
    for j in range(4) :
        if i == j :
            L[i][j] = 1
        elif i < j :
            L[i][j] = 2
            
print(L)
\end{verbatim}
\end{exampleblock}

\end{frame}

\begin{frame}[fragile]{Création de listes imbriquées}

Solution plus compacte: 


\begin{exampleblock}{}
\begin{verbatim}
L = [0] * 4
for i in range(4) : # pour chaque ligne
    L[i] = [0] * i + [1] + [2] * (n - i - 1)

print(L)
\end{verbatim}
\end{exampleblock}

\end{frame}
\end{comment}
\end{document}


\textbf{Remarque.} La fonction \texttt{len()} n'est pas propre aux listes
et peut s'appliquer à différents types d'objets.

\begin{exampleblock}{}
\begin{verbatim}
chaine = 'abracadabra'
print(len(chaine))
\end{verbatim}
\end{exampleblock}



\begin{frame}[fragile]{Chaînes de caractères}

\begin{exampleblock}{}
\begin{verbatim}
chaine = 'abracadabra'
chaine[2] = 'd'
print(chaine[2:5])
\end{verbatim}
\end{exampleblock}



\end{frame}



